% spring_math_summary.tex -- condensed main-body version of the spring
% control equations. Full derivations, the piecewise deadband, and the
% position-radius gate for pose springs are in the appendix
% (spring_math.tex); this fragment states only the equations actually
% run in the reported condition (no deadband, rest length zero).
%
% A fragment, not a standalone document: \input{} it from the paper.
% Requires: amsmath, amssymb, bm

\subsection{Virtual spring control}
\label{sec:spring-math-summary}

The arm is driven by a set of virtual springs, each mapping the current
configuration to a joint torque; torques are summed and commanded directly, so
no inverse kinematics or trajectory is involved. For a spring attached at a
point $\bm{p}(\bm{q})$ on link $\ell$, let $\bm{J}_{v}, \bm{J}_{\omega} \in
\mathbb{R}^{3\times n}$ be the point's translational and the link's rotational
Jacobian, so that a Cartesian force $\bm{f}$ and moment $\bm{m}$ map to joint
torque by the transpose relation
%
\begin{equation}
  \bm{\tau} \;=\; \bm{J}_{v}^{\top}\bm{f} \;+\; \bm{J}_{\omega}^{\top}\bm{m}.
  \label{eq:summary-transpose}
\end{equation}

\paragraph{Linear springs.} A linear spring pulls $\bm{p}(\bm{q})$ toward a
world target $\bm{p}^{*}$ with stiffness $k$ and damping $b$. Writing
$\bm{d}(\bm{q}) = \bm{p}^{*} - \bm{p}(\bm{q})$ for the displacement, the force
and resulting torque are
%
\begin{equation}
  \bm{\tau} \;=\; \bm{J}_{v}^{\top}\Bigl(\,\underbrace{k\,\bm{d}(\bm{q})}_{\text{elastic}}
  \;-\; \underbrace{b\,\bm{J}_{v}\dot{\bm{q}}}_{\text{damping}}\Bigr).
  \label{eq:summary-linear}
\end{equation}
%
(The full model additionally supports a dead zone and a soft-start ramp on
$\lVert\bm{d}\rVert$; see appendix.)

\paragraph{Pose springs.} A pose spring orients a link-fixed direction
$\bm{n}(\bm{q})$ toward a target heading $\bm{u}(\bm{q})$ (a look-at
constraint), rather than positioning a point. With axis--angle error
$(\theta, \hat{\bm{a}})$ between $\bm{n}$ and $\bm{u}$, rotational stiffness
$k_{\theta}$, and damping $b_{\theta}$,
%
\begin{equation}
  \bm{\tau}_{\mathrm{orient}} \;=\; \bm{J}_{\omega}^{\top}
  \Bigl(k_{\theta}\,\theta\,\hat{\bm{a}} \;-\; b_{\theta}\,\bm{J}_{\omega}\dot{\bm{q}}\Bigr).
  \label{eq:summary-pose}
\end{equation}
%
(The full model additionally gates the component of this torque that would
move $\bm{p}(\bm{q})$ once it drifts past a bounded radius; see appendix.)

\paragraph{Joint springs.} A small number of joints instead carry a torsional
spring acting directly on their own coordinate, with no Jacobian involved
--- needed for axes a Cartesian spring cannot see, e.g.\ a wrist joint whose
rotation does not move its own attachment point. For joint $j$,
%
\begin{equation}
  \tau_{j} \;=\; -\,k\,\Delta\theta \;-\; b\,\dot{\theta}_{j},
  \qquad
  \Delta\theta = \operatorname{wrap}_{(-\pi,\pi]}(\theta_{j} - \theta^{*}).
  \label{eq:summary-joint}
\end{equation}

\paragraph{Aggregation.} All springs superpose. With gravity compensation
$\bm{\tau}_{g}(\bm{q})$ and a ramp factor $s \in [0,1]$ applied only to the
spring sum (never to gravity compensation),
%
\begin{equation}
  \bm{\tau}_{\mathrm{cmd}} \;=\; s\sum_{i} \bm{\tau}_{i}(\bm{q}, \dot{\bm{q}})
  \;+\; \bm{\tau}_{g}(\bm{q}).
  \label{eq:summary-total}
\end{equation}
